\documentclass[11pt]{article}

\usepackage[margin=0.75cm]{geometry}
\usepackage{amsmath}
\usepackage{amssymb}
\usepackage{array}
\usepackage{enumitem}
\usepackage{graphicx}
\usepackage{svg}
\svgsetup{inkscapelatex=false}

\setlength{\columnsep}{0.55cm}
\setlength{\parindent}{0pt}

\newcommand{\hwsvg}[2]{%
  \IfFileExists{#1.svg}%
    {\includesvg[width=#2]{#1}}%
    {\fbox{\parbox[c][2.0cm][c]{#2}{\centering\scriptsize
      Run \texttt{python3 render\_homework4.py}\\to generate this figure.}}}%
}

\begin{document}

\twocolumn

\title{Homework 4, COSC330/PHYS306\\Chapters 5 and 6}
\author{Prof. Jordan C. Hanson}
\date{}
\maketitle

\section{Chapter 5}

\begin{enumerate}

\item \textbf{[Section 5-1, Exercise 6]}\\
Develop the truth table for each circuit in Figure 5--56.

\begin{center}
\hwsvg{fig_5_56a}{0.5\linewidth}\\[-0.05cm]
\textbf{(a)}

\vspace{0.12cm}
\hwsvg{fig_5_56b}{0.5\linewidth}\\[-0.05cm]
\textbf{(b)}

\vspace{0.12cm}
\hwsvg{fig_5_56c}{0.6\linewidth}\\[-0.05cm]
\textbf{(c)}

\vspace{0.12cm}
\hwsvg{fig_5_56d}{0.6\linewidth}\\[-0.05cm]
\textbf{(d)}

\vspace{0.12cm}
\hwsvg{fig_5_56e}{0.6\linewidth}\\[-0.05cm]
\textbf{(e)}

\vspace{0.12cm}
\hwsvg{fig_5_56f}{0.6\linewidth}\\[-0.05cm]
\textbf{(f)}

\vspace{0.10cm}
\textbf{FIGURE 5--56}
\end{center}

\item \textbf{[Section 5-2, Exercise 13]}\\
Implement a logic circuit for the truth table in Table 5--8.

\begin{center}
\renewcommand{\arraystretch}{1.10}
\begin{tabular}{c c c | c}
\hline
\multicolumn{3}{c|}{\textbf{Inputs}} & \textbf{Output} \\
$A$ & $B$ & $C$ & $X$ \\
\hline
0 & 0 & 0 & 1 \\
0 & 0 & 1 & 0 \\
0 & 1 & 0 & 1 \\
0 & 1 & 1 & 0 \\
1 & 0 & 0 & 1 \\
1 & 0 & 1 & 0 \\
1 & 1 & 0 & 1 \\
1 & 1 & 1 & 1 \\
\hline
\end{tabular}

\vspace{0.10cm}
\textbf{TABLE 5--8}
\end{center}

\item \textbf{[Section 5-3, Exercise 20]}\\
Implement the logic circuits in Figure 5--54 using only NAND gates.

\begin{center}
\begin{minipage}{0.48\linewidth}
\centering
\hwsvg{fig_5_54a}{0.99\linewidth}\\[-0.05cm]
\textbf{(a)}
\end{minipage}\hfill
\begin{minipage}{0.48\linewidth}
\centering
\hwsvg{fig_5_54b}{0.99\linewidth}\\[-0.05cm]
\textbf{(b)}
\end{minipage}

\vspace{0.10cm}
\textbf{FIGURE 5--54}
\end{center}

\end{enumerate}

\section{Chapter 6}

\begin{enumerate}
\setcounter{enumi}{3}

\item \textbf{[Section 6-2, Exercise 4]}\\
For the parallel adder in Figure 6--69, determine the complete sum by
analysis of the logical operation of the circuit. Verify your result by
longhand addition of the two input numbers.

\begin{center}
\hwsvg{fig_6_69}{0.96\linewidth}

\vspace{0.08cm}
\textbf{FIGURE 6--69}
\end{center}

\item \textbf{[Section 6-2, Exercise 8]}\\
The input waveforms in Figure 6--72 are applied to a 2-bit adder.
Determine the waveforms for the sum and the output carry in relation to
the inputs by constructing a timing diagram.

\begin{center}
\hwsvg{fig_6_72}{0.98\linewidth}

\vspace{0.08cm}
\textbf{FIGURE 6--72}
\end{center}

\item \textbf{[Section 6-4, Exercise 14]}\\
For the 4-bit comparator in Figure 6--74, plot each output waveform for
the inputs shown. The outputs are active-HIGH.

\begin{center}
\hwsvg{fig_6_74}{0.98\linewidth}

\vspace{0.08cm}
\textbf{FIGURE 6--74}
\end{center}

\item \textbf{[Section 6-5, Exercise 16]}\\
When a HIGH is on the output of each of the decoding gates in Figure
6--75, what is the binary code appearing on the inputs? The MSB is
$A_3$.

\begin{center}
\begin{minipage}{0.48\linewidth}
\centering
\hwsvg{fig_6_75a}{0.95\linewidth}\\[-0.05cm]
\textbf{(a)}
\end{minipage}\hfill
\begin{minipage}{0.48\linewidth}
\centering
\hwsvg{fig_6_75b}{0.95\linewidth}\\[-0.05cm]
\textbf{(b)}
\end{minipage}

\vspace{0.15cm}
\begin{minipage}{0.48\linewidth}
\centering
\hwsvg{fig_6_75c}{0.95\linewidth}\\[-0.05cm]
\textbf{(c)}
\end{minipage}\hfill
\begin{minipage}{0.48\linewidth}
\centering
\hwsvg{fig_6_75d}{0.95\linewidth}\\[-0.05cm]
\textbf{(d)}
\end{minipage}

\vspace{0.10cm}
\textbf{FIGURE 6--75}
\end{center}

\item \textbf{[Section 6-8, Exercise 28]}\\
For the multiplexer in Figure 6--79, determine the output for the
following input states:
\begin{equation*}
D_0=0,\quad D_1=1,\quad D_2=1,\quad D_3=0,
\end{equation*}
\begin{equation*}
S_0=1,\quad S_1=0.
\end{equation*}

\begin{center}
\hwsvg{fig_6_79}{0.88\linewidth}

\vspace{0.08cm}
\textbf{FIGURE 6--79}
\end{center}

\end{enumerate}

\end{document}
